% Generated draft: requires mathematical and source review.
\documentclass[11pt]{article}
\usepackage[margin=25mm]{geometry}
\usepackage{fontspec}
\setmainfont{Latin Modern Roman}
\usepackage{amsmath,amssymb,mathtools,unicode-math}
\setmathfont{Latin Modern Math}
\usepackage{enumitem,needspace,xurl,hyperref}
\setlength{\parindent}{0pt}
\setlength{\parskip}{5pt}
\setlength{\emergencystretch}{3em}
\newcommand{\N}{\mathbb{N}}
\newcommand{\Z}{\mathbb{Z}}
\newcommand{\Q}{\mathbb{Q}}
\newcommand{\R}{\mathbb{R}}
\newcommand{\C}{\mathbb{C}}
\newcommand{\F}{\mathbb{F}}
\newcommand{\A}{\mathbb{A}}
\newcommand{\PP}{\mathbb{P}}
\newcommand{\E}{\mathbb{E}}
\newcommand{\eps}{\varepsilon}
\newcommand{\dd}{\,\mathrm{d}}
\begin{document}
\section*{UnsolvedMath ID: 700554}
\textbf{Source problem code:} catalogue entry
\subsection{Definitions and mathematical statement}
For a nontrivial finite group \(X\), define its minimum proper-subgroup index by
\[
n(X)=\min\{[X:H]:H<X\},
\]
where \(H<X\) means that \(H\) is a proper subgroup of \(X\). This minimum exists because \(X\) is finite and nontrivial. The source does not specify a value for \(n(\{1\})\); the definitions below consequently use \(n\) only for nontrivial groups.

For a nontrivial finite group \(G\), let
\[
\mathcal S(G)=\{X<G:X\ne\{1\},\ n(X)=n(G)\}.
\]
A subgroup \(A\) is \emph{wide} in \(G\) if it is maximal under inclusion in \(\mathcal S(G)\). Explicitly, this means
\[
A\in\mathcal S(G)
\quad\text{and}\quad
\forall X\in\mathcal S(G),\quad A\le X\Longrightarrow A=X.
\]
Maximality here is within \(\mathcal S(G)\); the definition does not require \(A\) to be a maximal proper subgroup of \(G\).

\textbf{Part (a), retained in the main section according to the supplied record.} Determine all wide subgroups in finite projective special linear, symplectic, orthogonal, and unitary groups. These are the family names given in the source. The supplied record does not explicitly specify ranks, field sizes, simplicity restrictions, exceptional small parameters, or the precise projective and special forms intended for every family. Those scope questions remain unresolved; no additional exclusions or choices of forms are imposed here.

\textbf{Part (b), archived as solved according to the supplied record.} Prove, without using the classification of finite simple groups, that
\[
n(F_1)=[F_1:2F_2],
\]
where the source calls \(F_1\) and \(F_2\) the Fischer simple groups and describes \(2F_2\) as ``an extension of a group of order \(2\) by \(F_2\).'' The index requires \(2F_2\) to designate a subgroup of \(F_1\). The supplied record does not explicitly identify the Fischer groups in modern notation, specify that subgroup embedding, or disambiguate the extension description by an exact sequence. The displayed statement therefore retains the source notation and these unresolved qualifications. The supplied archive reports that this part was proved; this draft does not independently verify that report.
\subsection{Short English statement}
A wide subgroup is maximal among proper subgroups having the same minimum proper-subgroup index as the ambient group. Part (a) asks for their classification in the listed finite classical groups. Part (b) asks for a proof, without the classification of finite simple groups, of the stated minimum-index formula for Fischer groups. The supplied registry labels the combined problem partially solved and archives part (b) as solved. Family conventions and the notation in part (b) require review.
\subsection{Sources}
\begin{itemize}
\item Supplied local UnsolvedMath registry record: permanent ID \(700554\), source problem code \texttt{catalogue entry}, title \emph{Kourovka Notebook Problem 9.19}, upstream revision \texttt{d371ee917ac9647fbb06cd6fc529deec6f3dda82}. The supplied status is \texttt{partially\_solved}.
\item \emph{The Kourovka Notebook (Unsolved Problems in Group Theory)}, problem 9.19, proposed by V. D. Mazurov in the 9th issue (1984); the supplied background describes an update in the 21st issue (2026), identified as \texttt{arXiv:1401.0300v48}. The supplied collection description names E. I. Khukhro and V. D. Mazurov as editors of Issues 1--21 (1965--2026).
\item Supplied Archive of Solved Problems attribution for part (b): S. V. Zharov and V. D. Mazurov, in \emph{Matematicheskoye programmirovaniye i prilozheniya}, Ekaterinburg, 1995, pp. 96--97 (Russian). The source reports that the assertion was proved. No article title or fuller bibliographic identification is supplied.
\item The supplied status review is dated \texttt{2026-10-16}, labels the problem partially solved, and states that solved parts occur in the archive while remaining parts occur in the main section. This is supplied dated curation, not an independently verified literature assessment. No source was independently visited or checked for this draft.
\end{itemize}
\end{document}
